% Remove 12pt, change draftcls -> final when doing final.
\documentclass[journal]{./IEEE/IEEEtran}
\usepackage{cite,graphicx}

\newcommand{\SPTITLE}{Room-8! - A Matchmaking Application for Finding College Roommates}
\newcommand{\ADVISEE}{Joseph Matthew R. Marcos}
\newcommand{\ADVISER}{Caroline Natalie M. Peralta}

\newcommand{\BSCS}{Bachelor of Science in Computer Science }
\newcommand{\ICS}{Institute of Computer Science}
\newcommand{\UPLB}{University of the Philippines Los Ba\~{n}os }
\newcommand{\REMARK}{\thanks{Presented to the Faculty of the \ICS, \UPLB\
                            in partial fulfillment of the requirements
                            for the Degree of \BSCS }}

\markboth{CMSC 190 Special Problem, \ICS}{}
\title{\SPTITLE}
\author{\ADVISEE~and~\ADVISER%
\REMARK
}
\pubid{\copyright~2016~ICS \UPLB}

%%%%%%%%%%%%%%%%%%%%%%%%%%%%%%%%%%%%%%%%%%%%%%%%%%%%%%%%%%%%%%%%%%%%%%%%%%

\begin{document}

% TITLE
\maketitle

% ABSTRACT
\begin{abstract}
This article explores the role and significance of roommates in a college student's development - from the initial adaptation with the roommate setting, to his/her academic performance and life after college. We highlight the importance of finding the proper roommates so that a student can maximize his/her psychological development, get higher grades, and develop open-mindedness. We also propose a solution in the form of a web application that can help \UPLB students find mutually beneficial roommates for their college life. This solution also extends itself to solve the dorm-finding problem students experience at the start of every semester.
\end{abstract}

% INDEX TERMS
\begin{keywords}
UPLB, University of the Philippines Los Ba\~{n}os, Roommate, Roommate matching
\end{keywords}

% INTRODUCTION
\section{Introduction}
% To be effective, the introduction should answer the questions ``Why and What For (Four)?" Expanded, these questions are----
% \subsection{Why is the topic of interest?}

% Expound more on the psychological importance of college
Several studies support that social functioning during college plays a major role in a student's psychological development. According to Chickering and Reisser, one of the seven vectors of psychological developmental issues that college students face is cultivating mature interpersonal relationships\cite{chickering}. College-aged students also fall within Arnett's emerging adulthood stage, where identity formation is associated with friendships\cite{erb}. According to Erikson's stage theory of psychological development, a young adult prefers to seek intimacy in relationships rather than in isolation\cite{erikson}.

Studies have linked social functioning to mental health and adjustment to college life \cite{erb}. A plethora of benefits such as psychological well-being, environmental mastery, personal growth, and self-acceptance have all been related to a student's ability to form meaningful relationships\cite{erb}. The conjunction between a person’s social capabilities and these benefits is so strong that we can use it to predict how well a student  will adjust to the new demands in college life. An individual who can develop good relationships in college has less problems.
    \subsection{Why College Roommate Relationships}
    Our goal is to make it easier for students to acquire these benefits with less guess work. We want students to develop better interpersonal relationships in college $-$ particularly roommate relationships because these are relationships that are widely experienced by many college students\cite{erb}. Roommates have frequent contact, negotiations of responsibilities and compromises about the living environment, which often include noise level, cleanliness level, sleep/waking hours, visitors, decorations, and bills\cite{erb}. These may become challenges when dealing with each other.

    For many college students, many {"firsts"} add to the challenges of campus life. Usually, roommates are the first {"people of equal status"} they live with, compared to the parent-child relationship they usually have experienced at home\cite{erb}.

    Unlike friendships, college students may often not choose roommates, causing personality mismatches. Many studies support that roommates, while having the potential to be a boon to a student, can also be a hindrance to having a quality college life. According to Erb\cite{erb}, citing Liu\cite{liu}, in a 31,500 sample sized survey in the United States of America, 50.1\% of women and 44.1\% of men experienced “frequent” conflicts with roommates or housemates. In another survey in the same country, 5.6\% of undergraduates reported that difficulties with their roommates have hampered their academic endeavors\cite{erb}. For comparison, 4\% of American students said that alcohol did the same thing\cite{erb}.

    We target college roommates because it has been proven that college roommate relationships have a momentous effect on a student's performance before and after graduation. Although the college-roommate phenomenon has existed for a long time, there is only one paper that congregates empirical information about college-roommate relationships\cite{erb}.

    \subsection{Proposed Solution}
    In order to maximize the number of compatible roommates on campus, we will be making a roommate-matching application for UPLB. Unlike other roommate matching applications that exist today, we will target students of \UPLB that are looking for roommates. Students of \UPLB  have different needs and preferences to working class people, so we cannot use existing roommate matching applications. In addition, there are no roommate matching applications that cater to the UPLB area or the Philippines. It is about time we have one for campus.
\pubidadjcol
% Similar Applications
\section{Similar Applications}
There are existing roommate finding tools, which have more similarities than differences. They ask the user if they are looking for a place to stay, or if they have a room and are looking for someone to stay with them. Then, the user is \textit{bombarded} with a long questionnaire that s/he has to fill out. It covers most of the basic criteria you would have when looking for a room. The union of questions asked in both websites can be seen in the appendices.

% Both applications have similar questions, but have subtle insignificant differences. Roommates.com requires you to log in even before you start filling out the questionnaire. Roommie Match does not. However, in order to get a match, you will have to provide your credentials. Some people may find this surreptitious, as users will be forced togive their credentials after they have already filled out a long form.
\subsection{Roommates.com}
Roommate.com\cite{roommates.com} is a web application that caters to people looking for roommates in the United States of America. It can be accessed by going to \textit{www.roommates.com}. The home page enumerates cities where the application most often used, such as Las Vegas, Houston, Atlanta, Los Angeles, New York, etc....

Users are required to create an account to use the service. Account creation is free regardless of your location. Among its listed features for free accounts include: Photo profile, 2-way matching, and the freedom to contact potential roommates. The application has a premium account option, which is required to read the messages sent to you by other users in the application. It costs \$5.99 for three days, roughly \$20 dollars for 30 days, and \$30 dollars for 60 days.

Another feature of Roommates.com is the powersearch feature where the user can do searches with focused criteria. It is similar to Google's advanced search. The user will be redirected to a form that asks some questions (see appendix). Once the user has filled out the form, a list of matching roommates will appear. The user is able to see the profiles of the matched individuals, matching the information they wrote when creating an account.

The only way to contact a matched person is by sending them a message with the built in messaging system. Sending messages is free, but reading them requires a premium account. Therefore, you need to be a paying user to fully utilize Roommates.com.

Roommates.com is not feasible for use in UPLB because there are no locations near the \UPLB that come up in the list of city selections. The form also forces the user to specify his state - a piece of information that is irrelevant for people living outside the United States of America. The website has a lot of javascript alerts and prompts when navigating, which is an anti-pattern to most modern web applications (as of this writing in 2016). Finally, a free user won't be able to get any utility out of it without paying.

\subsection{Roommiematch.com}
Roommiematch is another roommate match making web application that caters to some cities in Canada and the United States of America. Unlike Roommates.com, Roommiematch starts asking users about their preferences without requiring them to create an account. It also caters to fewer cities than Roommates.com. It is more informal in its questionnaire - you can find statements like \textit{"People under 30 tend to party party PARTY!"} and \textit{"I affectionately keep dustbunnies as pets"}.

One of the strengths of Roommiematch.com over other web applications is that the forms are reviewed by actual humans before the pairing actually occurs. They claim that many of the registration forms get \textit{tossed}, ensuring the legitimacy of its subscribers. Roommiematch also has a more modern (2016) user interface - it has a flat design, and elements are larger than those found in Roommates.com.

Its weakness is that the registration process is so long and you cannot save your previous answers. It was composed of over ten multiple choice questions, and four questions where you had to type the answer. You have to answer all of them in one go.

A feature that may both be good and bad is that here is no way to view the profiles of people you are matched with except through the emails they send you. The reasoning behind this is so that less of the work goes to the user and their users' information does not get published everywhere. There is no need to revisit the website time after time. They claim that this is the feature that detters scammers from even registering.

Roomiematch can actually be considered a service that you subscribe to rather than an application. This is because there is no application you can interact with after filling out the form. Their team will email you updates every time you have a match, and it is up to you to contact your potential roommate.

\subsection{Tinder}
Tinder is a mobile application in Android and iOS we have to review as well. It's purpose is not to match potential roommates, but to match potential dating partners. We find it interesting because in 2015, it has an estimated 50 million users\cite{tinderstat2}, 80\% of which were were millenials, or individuals born after 1982 \cite{tinderstat}\cite{millenial}. It is intuitive to use, and has a use pattern that has been proven to be easily understood by millenials. We want our roommate matching application to have a similar use pattern to make it easier for the user.

When using Tinder, a user has to create an account. Tinder has to option of letting you log in with your Facebook account. Allowing Tinder to connect to your Facebook will allow it to get some data for your profile. There is an option to create an account without connecting it to Facebook. Once the account has been created, a user is prompted to fill out a few more fields, including a short description about oneself and a profile photo. The user can also make constraints that the application will consider when finding a match.

Tinder is all about the pictures and description a user entered. The user is shown a list of other users that may interest him/her. For every entry, the user can swipe right if s/he is interested in being matched with the user, and left if not. If two people are interested in each other, Tinder matches them together and the two users can communicate.

\section{The Ideal Roommate Matching application}
If we were to make an application, we want it to be similar to existing applications but tailor it in such a way that it caters to the interests of UPLB students. It has to be intuitive, easy to use, clean, and consistent. Finally, it has to have a \textit{no strings} attached approach, meaning we do not want to trap the user into giving us their credit card number out of frustration. This means we have to follow modern user experience strategies and plan out its features well to make it more relevant than the previous applications mentioned.

% MATERIALS AND METHODS
\section{Materials and Methods}
% Irving's Algorithm
\subsection{Irving's Algorithm}
Irving's Algorithm is an algorithm used to match two elements of the same set together. It solves what is commonly known as the Stable Roommates problem.

% Gale-Shapley's Algorithm
\subsection{Gale-Shapley's Algorithm}
Gale-Shapeley's Algorithm is an algorithm used to match two elements from different sets. It solves what is commonly known as the Stable Marriage problem.
% \section{Proof of methods down here}

% RESULTS AND DISCUSSION
% \section{Results and Discussion}

% CONCLUSION AND FUTURE WORK
% \section{Conclusion and Future Work}

\newpage
% APPENDICES
\appendices
% Please look at the IEEE guidelines
\section{Roommate.com and Roommiematch.com questions} %Mandatory title to specify an appendix
\begin{enumerate}
    \item Are you looking for a room?
    \begin{itemize}
        \item I am looking for a room
        \item I have a room available for rent
        \begin{itemize}
            \item Current household composition (How many people currently living there)
            \item Furnishing of dwelling
            \item Parking availability
            \item Does the room have its own bathroom
            \item Additional amenities (See to match to questions enumerated below)
        \end{itemize}
    \end{itemize}
    \item Who you are looking for
    \begin{itemize}
        \item Straight female
        \item Professional
        \item Non Smoker
        \item Lesbian
        \item Student
        \item Smoker
        \item Straight male
        \item Unemployed
        \item Outside Smoker
        \item Gay male
        \item Military
        \item Retired
    \end{itemize}
    \item Age range
    \item Who
    \begin{itemize}
        \item Needs a place to live
        \item Has a place available
    \end{itemize}
    \item Location
    \begin{itemize}
        \item State
        \item City
    \end{itemize}
    \item Available (Date the room is available or the roomate is available to move)
    \item Payment (price range)
    \item Cleanliness
    \begin{itemize}
        \item Average
        \item Clean
        \item Messy
    \end{itemize}
    \item Pets
    \begin{itemize}
        \item No dogs
        \item No cats
        \item No caged pets
    \end{itemize}
    \item Whether or not you have a preference regarding children in the home
\end{enumerate}

% Room-8! Questions - To be updated more when we have more research on optimal roommate traits
\section{Questions and their Category in Room-8!}
\begin{enumerate}

    \item What are you looking for?
    \begin{itemize}
        \item A place with vacancies
        \item I have a room with vacancies
    \end{itemize}

    \item When
    \begin{itemize}
        \item Summers only?
        \item Start date?
        \item Duration of stay?
    \end{itemize}

    \item Utilities
    \begin{itemize}
        \item Air conditioning?
        \item Can I do my own laundry?
        \begin{itemize}
            \item Is there a washing machine?
            \item Is there a dryer?
            \item Is there a place to dry my clothes?
        \end{itemize}
        \item Can I cook? (Check all that apply)
        \begin{itemize}
            \item Gas Stove
            \item Electric Stove
            \item Microwave oven
            \item Water kettle
        \end{itemize}
        \item Internet connection
        \begin{itemize}
            \item What is the minimum speed requirement?
            \item Do you torrent?
        \end{itemize}
    \end{itemize}

    \item Hobbies
    \begin{itemize}
        \item Are you a gamer?
        \item Are you an activist?
        \item Are you into sports
    \end{itemize}

    \item Lifestyle
    \begin{itemize}
        \item Are you religious? / Are you a part of a cult?
        \item Do you urinate on the shower?
        \item What is an acceptable noise level for you? (1 - 10)
        \item Do you drink Alcohol?
        \item Do you smoke?
        \item Do you prefer to study in the early mornings or evenings?
        \item Do you have guests?
        \item Are you willing to share a bathroom?
        \item Do you have guests come over?
        \item Do you live with pets?
        \begin{itemize}
            \item Small caged animals (birds, hamsters)
            \item Not small caged animals (dogs)
        \end{itemize}
        \item Are you willing to share a bathroom?
    \end{itemize}

    \item Location Preferences
    \begin{itemize}
        \item Nearby Restaurants? (5 minute walk)
        \item How far from UPLB (How long a walk)
        \item General Area (List Barangays)
    \end{itemize}

    \item Misc
    \begin{itemize}
        \item Study area for guests?
        \item Easily Accessible
        \item Cost
        \item Payment scheme
        \item Curfew
    \end{itemize}
\end{enumerate}


\newpage
% ACKNOWLEDGMENT
\section*{Acknowledgment}
I would like to thank my adviser \ADVISER for her patience.

% BIOGRAPHY
\begin{biography}[{\includegraphics{./yourPicture.eps}}]{\ADVISEE}
\end{biography}


\newpage
% BIBLIOGRAPHY
% \bibliography{./cs190-ieee}
\bibliographystyle{./IEEE/IEEEtran}
\bibliography{marcoscmsc190sp}
% \nocite{*}
\end{document}
